\documentclass[border=7bp]{standalone}
% Shared preamble for the report's editable vector figures.
\usepackage{fontspec}
\usepackage{xeCJK}
\setmainfont{Calibri}
\setCJKmainfont{SimSun}[AutoFakeBold=1.5]
\usepackage{amsmath}
\usepackage{tikz}
\usetikzlibrary{arrows.meta,positioning,calc,shapes.geometric,fit}
\definecolor{Ink}{HTML}{25313D}
\definecolor{Blue}{HTML}{315E80}
\definecolor{PaleBlue}{HTML}{E7EFF5}
\definecolor{Muted}{HTML}{626B74}
\definecolor{PaleGray}{HTML}{F0F1F3}
\tikzset{
  >=Stealth,
  every node/.style={font=\fontsize{10bp}{13bp}\selectfont,text=Ink},
  flow/.style={->,draw=Ink,line width=.65pt},
  frame/.style={rectangle,draw=Ink,line width=.65pt,align=center,
    minimum height=.85cm,text width=5cm,inner sep=5pt},
  emphasis/.style={frame,draw=Blue,fill=PaleBlue},
  supporting/.style={frame,draw=Muted,fill=PaleGray},
  annotation/.style={font=\fontsize{9bp}{12bp}\selectfont,text=Muted,align=center},
}

\begin{document}
\begin{tikzpicture}
\path (-7.3,0) rectangle (7.3,0);

\node[emphasis,text width=8cm] (read) at (0,0) {读取解路径 TSV：SolutionTrace};
\node[frame,text width=8cm] (validate) at (0,-1.4) {检查格式、元数据、步号、合法动作和最终目标};
\node[supporting,text width=3cm] (error) at (6.4,-1.4) {失败：抛出异常\\不创建播放窗口};
\node[frame,text width=8cm] (window) at (0,-2.8) {校验通过：创建 PuzzleVisualizer 窗口\\建立 PlaybackModel，当前步号置 0};
\node[emphasis,text width=8cm] (event) at (0,-4.3) {等待交互事件};
\node[frame,text width=4.2cm] (manual) at (-2.8,-6) {Previous / Next / Restart\\移动一步或回到初态，处理首末边界};
\node[frame,text width=4.2cm] (auto) at (2.8,-6) {Play / Pause / Speed\\控制 Swing Timer 与播放间隔};
\node[frame,text width=8cm] (refresh) at (0,-7.8) {更新当前步号、动作与棋盘；目标棋盘保持不变};
\draw[flow](read)--(validate);\draw[flow](validate)--(error);
\draw[flow](validate)--(window);\draw[flow](window)--(event);
\draw[flow](event.south)-|(manual.north);\draw[flow](event.south)-|(auto.north);
\coordinate (join) at (0,-7.1);
\draw[draw=Ink,line width=.65pt](manual.south)|-(join);
\draw[draw=Ink,line width=.65pt](auto.south)|-(join);\draw[flow](join)--(refresh.north);
\draw[flow](refresh.west)--++(-1,0)|-(event.west);
\node[annotation,text width=9cm] at (0,-8.85) {播放器从路径文件逐步还原棋盘；按钮与定时器共同控制进度，步号始终处于有效范围内};

\end{tikzpicture}
\end{document}
